\mwnewpart
\renewcommand\runningtitle{Section 3. Using RetrOS/LT}
\chapter{Using RetrOS/LT}
\sectioncontents

\section{The Prompt}\ixs{prompt}

When RetrOS/LT is ready for a command it shows a prompt, such as
\shown{A>}. The letter is the \term{current drive}, the one RetrOS/LT uses
when you do not name another. You type a command after the prompt and
press \key{Enter}.

\begin{dialogue}
A>\textbf{DIR}
\end{dialogue}

You may type commands in capitals or in small letters. If you make a
mistake before you press \key{Enter}, press \key{Backspace} to erase it. To
stop a command that is running, press \keys{Ctrl}{Break} or \keys{Ctrl}{C}.
To make the screen stop scrolling, press \keys{Ctrl}{Num Lock}; any key
starts it again.

To change the current drive, type its letter and a colon:

\begin{dialogue}
A>\textbf{B:}
B>
\end{dialogue}

After Install, the prompt also shows the current directory, as
\shown{C:\bs>}; the \typed{PROMPT} command sets what it shows.

\section{Drives}\ixs{drives}

Each diskette drive and each fixed disk has a letter. A and B are the
diskette drives; a computer with one diskette drive uses it as both A and
B, and asks you to change diskettes. C is the DOS partition of the first
fixed disk, and D that of the second, if there is one.

\section{Files and Their Names}\ixs{files}\ixs{file names}

Everything kept on a disk is in a \term{file}. A file has a name of one to
eight characters and, usually, an extension of up to three, with a period
between: \file{LETTER.TXT}, \file{BUDGET83.CAL}. You may use letters,
digits and these characters:

\begin{dialogue}
\$ \&\ \# @ ! \% ' ( ) - \{ \} \_ ` \textasciitilde\ \textasciicircum
\end{dialogue}

The extension usually tells what kind of file it is:

\begin{tabularx}{\textwidth}{@{}P{1.1in}X@{}}
\toprule
\file{.COM}, \file{.EXE} & Programs. Type the name, without the extension, to run one.\\
\file{.BAT} & Batch files: lists of commands, run the same way.\\
\file{.SYS} & Device drivers.\\
\bottomrule
\end{tabularx}

These names are devices, not files, and cannot be used as file names:
\file{CON} (the keyboard and screen), \file{PRN} and \file{LPT1} to
\file{LPT3} (printers), \file{AUX} and \file{COM1} (the serial port), and
\file{NUL} (nothing).

\subsection{Wildcards}\ixs{wildcards} Two characters stand for others in a
file name, so that one command can work on many files. A question mark
\typed{?} stands for any one character; an asterisk \typed{*} stands for
any characters to the end of the name or of the extension.

\begin{tabularx}{\textwidth}{@{}P{1.3in}X@{}}
\toprule
\typed{*.TXT} & every file with the extension \file{.TXT}\\
\typed{LETTER.*} & every file named \file{LETTER}, whatever its extension\\
\typed{CHAP?.DOC} & \file{CHAP1.DOC}, \file{CHAP2.DOC} and so on\\
\typed{*.*} & every file\\
\bottomrule
\end{tabularx}

\section{Directories}\ixs{directories}\ixs{path}

A fixed disk holds thousands of files. To keep them in order, a disk is
divided into \term{directories}, and a directory can hold other
directories as well as files. Every disk has one directory to begin with,
the \term{root}, written as a backslash: \file{C:\bs}.

The \term{path} to a file names the drive, then each directory on the way,
then the file, with backslashes between:

\begin{dialogue}
C:\bs{}LETTERS\bs{}1983\bs{}SMITH.TXT
\end{dialogue}

You may leave out what RetrOS/LT already knows. If the current drive is C
and the current directory is \file{\bs{}LETTERS\bs{}1983}, then
\file{SMITH.TXT} is enough. Each drive has its own current directory.

\begin{tabularx}{\textwidth}{@{}P{1.6in}X@{}}
\toprule
\typed{MD LETTERS} & makes a directory named \file{LETTERS} in the current one\\
\typed{CD LETTERS} & makes it the current directory\\
\typed{CD ..} & goes back to the directory above\\
\typed{CD \bs} & goes to the root\\
\typed{RD LETTERS} & removes the directory, which must be empty\\
\typed{DIR} & lists the files in the current directory\\
\typed{TREE C:} & shows all the directories of drive C\\
\bottomrule
\end{tabularx}

In every directory but the root, \typed{DIR} shows two entries named
\file{.} and \file{..}; they stand for the directory itself and the one
above it. The root directory of a disk holds a fixed number of entries:
64 on a single-sided diskette, 112 on a double-sided one, and 512 on a
fixed disk. Other directories can hold any number.

\section{Working with Files}

\begin{tabularx}{\textwidth}{@{}P{2.2in}X@{}}
\toprule
\typed{COPY A:REPORT.TXT C:} & copies a file to another disk\\
\typed{COPY REPORT.TXT REPORT.BAK} & makes a second copy under another name\\
\typed{COPY A:*.* C:} & copies every file of a directory\\
\typed{REN REPORT.TXT FINAL.TXT} & changes a file's name\\
\typed{DEL REPORT.BAK} & erases a file\\
\typed{TYPE REPORT.TXT} & shows a file on the screen\\
\typed{PRINT REPORT.TXT} & prints it while you go on working\\
\bottomrule
\end{tabularx}

\caution{\typed{DEL} and \typed{COPY} do not ask before they erase or
replace a file, and RetrOS/LT has no command to bring an erased file back.
Only \typed{DEL *.*} asks first.}

\section{Running Programs}\ixs{programs, running}

To run a program, type its name. RetrOS/LT looks for it in the current
directory, and then in each directory named by the \typed{PATH} command.
Install sets the path to the directory RetrOS/LT is in, so its commands
can be used from anywhere:

\begin{dialogue}
C:\bs>\textbf{PATH C:\bs{}DOS;C:\bs{}WP}
\end{dialogue}

Some commands, such as \typed{DIR} and \typed{COPY}, are part of
\file{COMMAND.COM} and are always there; they are called \term{internal}.
The rest are programs on the disk, and are called \term{external}.
Section~6 says which is which.

\section{Redirection and Pipes}\ixs{redirection}\ixs{pipes}

Most commands write to the screen and read from the keyboard. You can send
what a command writes to a file or a device, or have it read from a file:

\begin{tabularx}{\textwidth}{@{}P{2.2in}X@{}}
\toprule
\typed{DIR > LIST.TXT} & writes the listing to a file, replacing it\\
\typed{DIR >> LIST.TXT} & adds the listing to the end of the file\\
\typed{DIR > PRN} & prints the listing\\
\typed{SORT < NAMES.TXT} & sorts the lines of a file\\
\typed{DIR | SORT} & sends the listing through \typed{SORT}\\
\typed{TYPE BIG.TXT | MORE} & shows a file a screen at a time\\
\bottomrule
\end{tabularx}

The bar \typed{|} is a \term{pipe}: it makes what one command writes into
what the next one reads. RetrOS/LT passes it through a temporary file in
the root directory of the current drive, so that drive must not be
write-protected. \typed{SORT}, \typed{FIND} and \typed{MORE} are made to be
used this way.

\section{The Keyboard}\ixs{editing keys}

RetrOS/LT remembers the last line you typed. These keys bring it back,
whole or in part:

\begin{tabularx}{\textwidth}{@{}P{1.1in}X@{}}
\toprule
\key{F1} or \key{$\rightarrow$} & copies one character from the last line\\
\key{F2} \var{c} & copies up to the character \var{c}\\
\key{F3} & copies all the rest of the last line\\
\key{F4} \var{c} & skips up to the character \var{c}\\
\key{F5} & keeps the line you have typed for more editing, without running it\\
\key{Ins} & inserts what you type; \key{Del} skips a character\\
\key{Esc} & cancels the line\\
\bottomrule
\end{tabularx}

Other keys work at any time:

\begin{tabularx}{\textwidth}{@{}P{1.6in}X@{}}
\toprule
\keys{Ctrl}{Break}, \keys{Ctrl}{C} & stops the command that is running\\
\keys{Ctrl}{Num Lock} & stops the screen from scrolling, until any key is pressed\\
\keys{Shift}{PrtSc} & prints what is on the screen\\
\keys{Ctrl}{PrtSc} & starts printing everything that appears on the screen; press again to stop\\
\key{Ctrl}\,+\,\key{Alt}\,+\,\key{Del} & restarts the computer\\
\bottomrule
\end{tabularx}

\section{Batch Files}\ixs{batch files}

A batch file is a file of commands, one to a line, with the extension
\file{.BAT}. Type its name and RetrOS/LT runs the commands in turn. You
can make one with \typed{EDLIN}, or by copying from the keyboard:

\begin{dialogue}
A>\textbf{COPY CON SAVE.BAT}
\textbf{ECHO OFF}
\textbf{ECHO Put the backup diskette in drive B}
\textbf{PAUSE}
\textbf{COPY C:\bs{}WORK\bs{}*.* B:}
\key{F6} \key{Enter}
\end{dialogue}

Within a batch file, \typed{\%0} stands for its own name and \typed{\%1} to
\typed{\%9} for the words typed after it. The commands made for batch
files are \typed{ECHO}, \typed{FOR}, \typed{GOTO}, \typed{IF},
\typed{PAUSE}, \typed{REM} and \typed{SHIFT}; Section~6 describes each.
\typed{ECHO OFF} stops the commands from being shown as they run, but is
itself shown.

The batch file \file{AUTOEXEC.BAT} in the root of the startup disk is run
each time the computer starts, in place of the date and time questions.
See Section~5.

\section{BASIC}\ixs{BASIC}\ixs{BASICA command}\ixs{GW-BASIC}

RetrOS/LT's BASIC is Microsoft's GW-BASIC, built from the source
Microsoft published, with RetrOS/LT's own part for the IBM Personal
Computer's screen, keyboard, graphics, speaker, printers and serial port.
It understands the BASIC of IBM's \typed{BASICA} for PC DOS 2.10, and
programs written for that run in it. It is \file{BASICA.EXE}, on the
Supplemental Programs diskette; \file{BASIC.COM} runs it, so
\typed{BASIC} and \typed{BASICA} both start it:

\begin{dialogue}
C:\bs>\textbf{BASICA}
RetrOS/LT BASIC Version 1.0
(C) Copyright Microsoft 1982
62746 Bytes free
Ok
\end{dialogue}

Type \typed{SYSTEM} to leave BASIC. \typed{BASICA CLOCK} starts BASIC and
runs the program \file{CLOCK.BAS}. IBM's switches are accepted; see
\typed{BASIC} in Section~6.

\file{BASIC.COM} looks for \file{BASICA.EXE} in the current directory and
then in the \typed{PATH}, as RetrOS/LT looks for a program. Once Install
has set the path, \typed{BASIC} works from anywhere; from the diskettes,
make the Supplemental Programs diskette's drive the current one first, or
type \typed{BASICA}.

BASIC takes about 53 KB for itself and up to 64 KB more for programs and
their data, all it can address; with less memory, it uses what is left.

\subsection{The sample programs} The Supplemental Programs diskette has
six, written for RetrOS/LT. Run one with \typed{BASICA} and its name, or
\typed{RUN "}\var{name}\typed{"} in BASIC; \typed{LIST} shows how it is
done.

\begin{tabularx}{\textwidth}{@{}P{1.1in}X@{}}
\toprule
\file{CLOCK.BAS} & A clock with hands, set by \typed{TIME\$}.\\
\file{TUNES.BAS} & Three old tunes, played with \typed{PLAY}.\\
\file{PALETTE.BAS} & The text colors, and the two graphics palettes.\\
\file{SKETCH.BAS} & Drawing with the arrow keys.\\
\file{LOAN.BAS} & A loan's monthly payment, and its first year.\\
\file{BOUNCE.BAS} & A ball, with \typed{GET}, \typed{PUT} and \typed{SOUND}.\\
\bottomrule
\end{tabularx}

The graphics programs need the Color/Graphics Adapter.

\subsection{What it does not have} It has no cassette device,
\file{CAS1:}, since it runs from disk. \typed{TIMER}, \typed{CHDIR},
\typed{MKDIR}, \typed{RMDIR}, \typed{SHELL}, \typed{ENVIRON},
\typed{IOCTL}, \typed{ERDEV}, \typed{VIEW}, \typed{WINDOW} and
\typed{PMAP}, which later versions of BASIC added, answer
\shown{Advanced Feature}.

\subsection{Cassette BASIC}\ixs{Cassette BASIC}\ixs{ROM BASIC} The IBM
Personal Computer, the XT and some later IBM computers have BASIC in
read-only memory: Cassette BASIC, the BASIC the computer starts when it
finds no disk to start from. \typed{BASIC /ROM} (or \typed{BASICA /ROM})
leaves RetrOS/LT and starts it:

\begin{dialogue}
C:\bs>\textbf{BASIC /ROM}
\end{dialogue}

On a computer without it, BASIC says \shown{No ROM BASIC in this computer}
and you are back at the prompt.

\caution{Cassette BASIC cannot use diskettes or the fixed disk. The only
place it can keep a program is a cassette recorder, plugged into the
cassette connector of the original IBM Personal Computer; the XT, the AT
and the PS/2 have no such connector, and on them what you type in
Cassette BASIC cannot be saved at all. There is no way back to RetrOS/LT
but restarting the computer, which loses whatever was in Cassette BASIC's
memory. Save your work before you leave.}

\section{Restarting}

To restart the computer without turning it off, press
\key{Ctrl}\,+\,\key{Alt}\,+\,\key{Del}. Wait until the drive lights are out
and the prompt is showing before you do, or before you turn the computer
off, so that everything has been written to the disk.
